\chapter{How the data was made: the engineering behind the figures}

Every figure and every quotation in this note comes out of a pipeline, and a reader who wants to argue with a number needs to know which stage produced it and what that stage refused to do. This chapter describes the pipeline as it stood on 13 September 2026, stage by stage, in the order the data flows. It is the edition's method, not mine, and where it is quoted the words are the scripts' own headers.

\section{The inventory: harvested by a person, committed as text}

The Whitney publishes its photographs of the ledgers through ResourceSpace, one image to a resource reference, with a descriptor for each (\q{Page 62 [horizontal orientation, ``Fall 1927\ldots'']}). The images answer a plain HTTP request; the HTML pages that list them sit behind a JavaScript bot-check, and the project does not impersonate a browser to get past it. So the inventory of the ten collections, six ledgers and four notebooks, was harvested by a person in a browser with a twelve-line snippet, and the result was committed as ten text files under \path{harvest/}. From those, \path{scripts/catalogue.mjs} generates a typed catalogue in which the museum's descriptors survive unaltered; the one thing it infers is the leaf number the Hoppers wrote on the paper, parsed out of the descriptor and marked as inferred. The continuous integration workflow regenerates the catalogue on every push and fails if the committed file differs, so the catalogue cannot drift from the harvest.

Three numbers overlap on every sheet and the data model keeps them apart: the leaf number written on the paper, which names six different things across six volumes; the museum's sequence number, which counts photographs, covers and versos included; and the ResourceSpace reference, which is in no order at all and is the only identifier that names exactly one image. Every citation in this note prints the leaf and the reference for that reason.

\section{Nothing of the archive is stored}

Not one photograph is in the repository or on any disk of the project's. The reading view puts an image element straight at the Whitney's file, fetched as it is looked at; the \texttt{.gitignore} refuses \texttt{*.jpg}. This was measured rather than assumed: on 5 September 2026 the museum's download endpoint was found to answer cross-origin requests, to echo the requesting origin into its CORS header, to hold a current certificate, to check no referer and to honour byte ranges, so no relay is needed. The largest published size, 1292 by 2000 pixels, is about 170 pixels per inch across a leaf, and was checked to be legible for Josephine Hopper's figures before anything else was built. Only a transcription pass needs sheets on disk; \texttt{npm run mirror} fetches twelve at a time into a git-ignored directory, sleeping ten seconds between sheets as the museum's \texttt{robots.txt} asks, and \texttt{npm run tiles} cuts each into overlapping crops at twice size, because \q{choosing crop rectangles is not reading}.

\section{The transcription: a restricted language, and what it refuses}

A batch is twelve sheets, read in one fresh context by a model the skill accepts (Fable 5.1, Fable 5 or Opus 5; any other is refused), and written as a LaTeX file whose header names the model, the date, and what was hard. The file uses a deliberately small set of macros, defined in \path{transcripts/preamble/hopper.sty}, and each one is a claim of a particular kind: \verb|\sheet{ref}{leaf}| opens a sheet and is what turns the photograph when the transcript is scrolled; \verb+\hand{}+ attributes a hand to Josephine, to Edward, to a later writer or to nobody, and \emph{unidentified} is a permitted and frequent answer; \verb|\ill{}| marks what could not be read and is never guessed; \verb|\uncertain{}| offers a reading and flags it; \verb|\struck{}| keeps what the writer crossed out; \verb|\note{}| is the transcriber's, \verb|\marginal{}| the Hoppers'; \verb|\sketch{}| stands where Edward Hopper's record drawing is; \verb+\ink{}+ records the colour, red, pencil or blue, of a cell where the leaf's colour means something; \verb|\entry{as written}{ISO date}| opens a diary entry with the writer's date and the transcriber's date kept apart; and \verb|ledgertable| carries the ruled columns.

The governing rule is stated in the method: \q{a guessed figure is a sale that did not happen}. Nothing is expanded, nothing is normalised, nothing is corrected; \q{Inv.} stays \q{Inv.}, \q{30 - 1/3} stays as written, and a name spelt four ways stays spelt four ways. There is one edition and no second register, because a ledger's English is plain and a second text would be a second artifact to keep in step. The renderer, \path{scripts/render.mjs}, accepts only that subset and fails loudly outside it, naming the file, the line and the construct, on the principle that \q{a silently mangled ledger is a table of numbers that look right}; it also refuses a resource reference that is not in the harvest and a file that omits the rights line.

What the transcription is not is checked. \path{transcripts/status.json} can hold two states no file can prove, \emph{checked} and \emph{skipped}, and it holds neither for any batch. Everything else the site says about progress is observed by \path{scripts/manifest.mjs} from the files on disk, \q{so that a progress figure can never claim work that was not done}. One reading in Book I is marked as checked by a person; the rest is a first pass.

\section{The derived files, and the rule that a build must fail}

From the \texttt{.tex} alone, scripts derive everything the site and this note use, and each script's header says what it counts and what it therefore is not.

\path{scripts/tei.mjs} writes a TEI P5 file per transcription by a one-to-one mapping of the macros onto TEI elements, \verb|\ill{}| to \texttt{<gap reason="illegible"/>}, \verb|\uncertain{}| to \texttt{<unclear>}, and re-reads no sheet; \path{scripts/tei-validate.mjs} validates the result against the project's own customisation of fourteen element specifications, because, as its header records, the export was well-formed XML and invalid TEI for as long as nothing had asked. \path{scripts/citation.mjs} generates the citation files from a census of what is transcribed, on the ground that a hand-written figure \q{is wrong from the next batch onward}. \path{scripts/glossary.mjs} builds the key to the abbreviations and price notations from hand-declared entries, each carrying a search pattern, and fails the build if an entry matches no leaf; the search runs over the books' own writing and never over a \verb|\note{}|, so that an entry cannot cite as evidence the sentence written because of it. \path{scripts/framing.mjs} applies the same rule to the black notebook's frames and sums. \path{scripts/streak.mjs} resolves every dated row to a calendar day and says, in its header, that the result records the days the books record something and not the days anybody worked. \path{scripts/formats.mjs} reads the size out of the title line only, in his hand, and rounds nothing. \path{scripts/materials.mjs} reads the paint formula out of \verb|\hand{}| bodies only, never out of a note or a marginal, and counts once per work. \path{scripts/works.mjs} retrieves the museum holdings of each titled work from five institutions' public APIs, checks the artist field on every record, and states that a link says a work of this title is held there and not that the row concerns that copy. \path{scripts/notes.mjs} refuses a per-work note whose source key does not resolve. \path{scripts/places.mjs} refuses a journey whose sourced event it cannot find in the timeline. The CI workflow runs the catalogue, the renderer, the TEI export and validation, the manifest, the citation, the glossary, the framing and the spending scripts, compares each generated file with the committed one, and only then builds the site.

\section{The accounts reader}

The figures in the accounts chapter and Tables A.1 to A.6 come from \path{scripts/accounts.mjs}, which has two readers because the books have two shapes.

The work-books reader counts a row only where the leaf puts a price in one cell and a date in another, because the column a date stands in is what says whether it dates the sale or the cheque. On the oils leaves the first column dates the sale; on the etchings leaves it dates the exhibition, so those rows are set aside as unparsed, 201 of them, rather than dated wrongly. Book III rules nothing and is read for the sale sentence, \q{6000 - 1/3 = 4000. Feb. 17}, and nothing else. Where a row states a price, a rate and a receipt, the receipt is checked against price less commission; a disagreement is reported and never corrected. A sale written into two volumes is matched on title, price and rate and counted once, and the collapse is published so it can be argued with. A party is credited on its name or surname; two spellings merge only where one's words sit inside the other's, the remaining merges are declared by hand in an aliases file, and a bare surname shared by two declared parties is reported ambiguous and credited to nobody. Money is kept in three buckets, collected, still owed and not recorded, and the script asserts at every level that they sum to the accrual.

The Book IV reader is the more consequential and its rules are set out in \path{docs/book-iv-method.md}, period by period, because the writer's method changes four times. Items are the charge; a \q{bill} counts only where nothing was itemised since the last settlement, or every commission in the volume would be counted twice. A bare figure is a subtotal where a bill, a \q{less} line or a carried total follows it, or where a cheque follows it and the priced lines above add to it; the leaf rules a line above a subtotal and the transcription does not record the rule, so the rule is recovered from the arithmetic. A \q{less} line is a deduction whichever thing it names. What stands between it and the cheque is the net restated. A red receipt rubbed out and rewritten is counted once. A title charged more than once, where one entry says on account or balance, is one picture at the largest figure and the others are reported as instalments. The year is carried from the last one the volume states, and it carries across a batch boundary only where the batches abut. Receipts are kept as their own stream, because a cheque settles a run of work and the runs overlap. The check on all of it is the book's own pencil totals, and the reader's receipts reconcile to them to the cent wherever the leaf's arithmetic closes.

Two of these rules were found by bugs. An unanchored match on the word \emph{bill} took a story title, \q{Wild Bill in Deadwood Gulch}, for an invoice and lost a drawing of 1915; a header rule read \q{May 25th Initial Letters} as a new client and counted an Adventure bill of 225 twice. Both were caught by checking the per-year totals against leaves read by hand, and the method note records that a faithful transcription can surface a reader bug without containing one.

\section{The ink pass}

Book IV's colour is its status field, and fourteen batches were transcribed without recording it. \path{scripts/ink.mjs} reads the colour off the photographs, because \q{the colour is the one thing on the leaf a machine can see better than the hand it is written in}: it aligns each leaf's transcribed rows with the written lines of its image, classifies each money cell as black, red or pencil, and writes the result into \verb|\ink{}| and \verb|\ruledoff| where the transcriber accepts it. Where it has run, twelve leaves so far, red money is a receipt and pencil money a sum before any word is read. It also found, on leaf 21, a whole entry the transcription lacked, and it reports seventy-four leaves with at least one written line the transcription does not have, which is the worklist for the re-reading pass.

That figure needs stating for what it is, because it bears on this note and not only on the accounts. The ink pass has run on Book IV alone, so the seventy-four leaves are Book IV leaves, about half of that volume; many of the lines it flags are the alignment's own noise, and each has to be looked at, but at least one, the Country Gentleman entry of 26 February 1917 on leaf 21, was a real omission. The consequence for this note is bounded and can be said exactly. The Book IV chapter quotes the pocket book for its shape and its last leaves, and those quotations are of lines the transcription has, not of lines it lacks; what a missing line changes is a sum, so the figures of the cash book, Tables A.1 and A.2 and Figure 10.1, are floors in a second sense as well as the first: a leaf may hold a charge or a receipt the reader never saw. No such pass has run on the other five volumes or the notebooks, so nothing is known about their omissions, which is a different and larger uncertainty than a counted one. The quotations from Books I to III and from the notebooks in this note are readings of lines that exist on the sheet; whether a neighbouring line was skipped is not something any reading here can rule out.

\section{The working method: a person, two models, a repository and an interface}

The pipeline above was not bought and was not assembled from a transcription service. It was built by one person, Michel Hua, over a fortnight in September 2026, and the way it was built is itself part of the method, because it is what makes the reading checkable.

\subsection*{Claude as the reader}

The transcriptions were made with Claude Code, Anthropic's command-line agent, running the models Claude Opus 5 and Claude Fable 5.1. The work is packaged as skills, small instruction files kept in the repository under \path{.claude/skills/}: \texttt{transcribe-hopper} reads a batch of twelve sheets into the restricted LaTeX and writes its header; \texttt{tag-hopper} writes the keywords line that closes a transcription; \texttt{modernize-hopper} writes the per-work notes from named museum sources only. The skills accept Fable 5.1, Fable 5 or Opus 5 and refuse any other model, and every file records which one read its sheets and on what day, so that a batch's provenance is a fact about the file and not about whichever model happened to be selected. Opus 5 read most of the ledgers between 5 and 12 September; Fable 5.1 read the last sittings of the black notebook and the Battle of Wash Sq. on 10 and 13 September, and drafted this note. The model does what a first transcriber does and no more: it reads what it can, marks what it cannot, attributes hands cautiously, and writes a note where a decision was made. It does not check itself, and the edition does not pretend it did.

\subsection*{GitHub as the workbench}

The repository on GitHub is the record of the work and the place it is done. Every batch, every script and every correction is a commit on the main branch, and the issue tracker holds the open questions as numbered issues with a plain convention, a \q{[P1]} in the title meaning that the issue blocks the edition: the ink of Book IV, the Zenodo deposit, the per-file revision history a citation will one day need. The continuous integration workflow described above runs on every push and refuses a push whose generated files have drifted from their sources, and the same workflow builds and deploys the site to GitHub Pages under the project's own domain. Nothing about the edition lives on a machine of its author's: a clone of the repository is the edition entire, minus the photographs, which are the museum's.

\subsection*{An interface built for the purpose}

The site at hopper.commutator.io was written by Michel Hua for this project, and it does one thing that no general-purpose tool did: it puts the transcription beside the Whitney's own photograph of the sheet, so that scrolling the reading turns the photograph and a doubtful figure can be checked against the hand it came out of on one screen. Everything else on it follows from that. The reading view offers a Cite button that builds the three citations, sheet, leaf and batch, from what the page already knows; the Accounts, Technique, Timeline and Schema tabs are the derived files rendered, with the rules that produced them printed beside the figures; the Diaries side serves the four notebooks the same way. The interface is the checking instrument. A machine reading that could not be put beside its source would be a text nobody could correct, and the whole value of the transcription is that it can be.

\subsection*{Working in parallel: Claude Code and remote control}

A corpus of five hundred sheets is not read in one sitting, and it was not read in one session. Claude Code runs as a local session on the author's machine, and several sessions can run at once: one transcribing a batch of Book IV while another reads the ink off its photographs, another writes the accounts reader and a fourth builds a page of the site to show the result. Each session has its own context, which is what the skill demands of a batch, and each leaves its work as a commit the others pull. Claude Code's remote control lets those sessions be started, directed and followed from another device, so that the machine at 3 Washington Square North's equivalent, a desk in one place, keeps working while its owner is elsewhere and answers a question from a telephone. That is an interesting way to work for a reason that has nothing to do with convenience: it makes the parallelism explicit. A transcription pass and an analysis pass over the same leaves are different jobs with different rules, and running them as different sessions, each with its own instructions and its own commit, keeps them from contaminating each other in the way a single long conversation would. This note was made the same way: eight readings of the transcripts ran at once, one per volume or notebook, and the prose was written from their reports.

\section{What this note added, and how}

This note was written in one session on 13 September 2026. Eight readings of the transcripts were run in parallel, one per volume or notebook, each instructed to read its files whole and return the structure, the quotable passages exactly as transcribed with leaf and reference, the hard data, the hands, and the transcriber's notes; the dealers' book was read in two passes when its first reading found only a quarter of the leaves. The tables in the appendix were generated by a small script from the edition's derived files and are reproducible from them. The prose was then written from those readings and from the derived data, chapter by chapter, and compiled with Tectonic. Nothing in it was checked against a photograph. Where the readings disagreed with a batch header, both are reported; where a reading found a defect in the transcription, it is listed for correction at source and not corrected here, on the edition's own rule that a corrected reading goes to the \texttt{.tex} with the name of who checked it.

That is the chain. The photograph is the museum's; the harvest is a person's; the transcription is a model's, first pass, dated; the derived files are scripts', with their rules published and their failures loud; the note is a model's again. Every link can be inspected, and the weakest one, the unchecked reading, is the one every sentence here says it rests on.
